\chapter{Metodología de Desarrollo}\label{chapter04}

% Detallar el proceso de construcción, las tecnologías utilizadas y las decisiones de diseño.

\section{Organización y herramientas de trabajo}
\label{sec:herramientas_metodologia}

El control de versiones del proyecto corre sobre GitHub, tanto para el código del sistema (Sección~\ref{sec:diagrama_componentes}) como para este mismo informe, escrito en LaTeX y versionado en un repositorio separado. Mantener el informe bajo control de versiones, en lugar de un editor de texto colaborativo convencional, permite llevar un historial de cambios explícito por sección y trabajar en paralelo sobre distintos capítulos sin sobrescribirse entre los dos integrantes.

Los diagramas presentados en el Capítulo~\ref{chapter03} (Diagrama de Casos de Uso, de Flujo de Información, de Componentes y de Arquitectura) se elaboraron en Lucidchart, en lugar de una herramienta de dibujo genérica, porque ofrece plantillas y notación estándar para diagramas UML y de arquitectura, lo que agiliza mantener una consistencia visual entre los distintos diagramas del informe.

Para el material que no se versiona bien como archivo de texto plano (las grabaciones y transcripciones de las entrevistas semiestructuradas y los datos de la encuesta del User Research, Capítulo~\ref{chapter03}) se utiliza Google Drive como repositorio compartido.

La planificación y el seguimiento de tareas del equipo se gestionan en Notion, donde se desglosan los hitos del cronograma en tareas concretas asignadas a cada integrante, mientras que el seguimiento con el tutor se realiza mediante reuniones periódicas por Microsoft Teams.

\section{Tecnologías}
\label{sec:tecnologias}

El Marco Teórico (Sección~\ref{sec:vlm_marco}, Sección~\ref{sec:edge_marco}) ya justificó, desde una perspectiva académica y clínica, por qué YOLO, CLIP, el clustering dinámico y el paradigma de edge computing son las técnicas adecuadas para el problema que aborda este proyecto, y el Diagrama de Componentes (Sección~\ref{sec:diagrama_componentes}) y el Diagrama de Arquitectura (Sección~\ref{sec:diagrama_arquitectura}) ya mostraron cómo está organizado el código validado hasta esta entrega y cómo se despliega el sistema completo, desde la cámara IP hasta la aplicación móvil. Sobre esa base conceptual y de despliegue, lo que sigue es la decisión de ingeniería inmediatamente debajo de esas dos capas, es decir, qué lenguaje usa cada componente, qué frameworks y librerías concretas lo implementan, y qué protocolo de comunicación lo conecta con el resto del sistema, además de por qué se eligió cada uno frente a sus alternativas.

\subsection{Lenguaje y stack de desarrollo}

Todo el sistema, tanto el pipeline de visión que corre en el borde como el backend que va a correr en la nube, está escrito en Python (versión 3.12 o superior). La decisión responde en primer lugar al ecosistema, dado que las librerías de referencia para cada una de las tareas del sistema (Ultralytics para pose estimation, ONNX Runtime para la inferencia de CLIP, scikit-learn para clustering) están escritas o tienen sus \textit{bindings} principales en Python, y reescribir el pipeline en un lenguaje de más bajo nivel, como C++, hubiese significado renunciar a esas librerías o reimplementarlas sin necesidad, dado que RNF4 exige una operación cercana al tiempo real y no tiempo real estricto. En segundo lugar, usar el mismo lenguaje en el borde y en la nube reduce la fricción de mantener dos bases de código con convenciones distintas, y permite, en particular, reutilizar sin cambios el esquema de eventos definido con Pydantic (Sección~\ref{sec:diagrama_componentes}) tanto en el almacenamiento local de \texttt{events/} como en el futuro backend que lo va a exponer sobre la nube. El costo de esta decisión es previsible, ya que Python introduce un overhead de interpretación que un lenguaje compilado no tendría, pero ese costo ya está absorbido en el diseño, dado que ni la estimación de pose ni la inferencia de CLIP corren sobre el intérprete de Python en sí, sino delegadas a Ultralytics y a ONNX Runtime respectivamente, ambas con su núcleo de cómputo implementado en C++ y expuesto a Python solo como interfaz.

\subsection{Frameworks y librerías por módulo}

\subsubsection*{Ultralytics, OpenCV y NumPy en \texttt{vision/}}
El módulo \texttt{vision/} implementa la heurística de detección de caídas e inmovilidad (Sección~\ref{sec:diagrama_componentes}) sobre keypoints de pose obtenidos con YOLO11n-pose a través de la librería Ultralytics. Implementar la arquitectura de YOLO desde cero, o recurrir a una implementación de terceros menos consolidada, hubiera significado resignar lo que Ultralytics ya resuelve, una interfaz de alto nivel que integra la descarga del modelo, la inferencia y la exportación a otros formatos. Además, es la misma línea de modelos que valida \textcite{Giri2025} y que confirma \textcite{Wang2025} como la variante de referencia para hardware embebido (Sección~\ref{sec:edge-inference}), de modo que usarla evita apartarse de la literatura que ya validó la viabilidad del enfoque. OpenCV se usa para todo el manejo de frames (lectura de archivos y de streams, recorte, dibujo del overlay de esqueleto) y NumPy para las operaciones vectoriales sobre los keypoints (velocidad vertical, ángulo de orientación) que alimentan la heurística de caída, ambas librerías estándar de facto en visión artificial con Python.

\subsubsection*{ONNX Runtime para la clasificación semántica}
CLIP se entrena y se exporta una única vez, en la máquina de desarrollo, con PyTorch, pero en el dispositivo de inferencia (hoy la PC de desarrollo, más adelante la Raspberry Pi) corre exclusivamente sobre ONNX Runtime, sin dependencia de PyTorch en tiempo de ejecución. Esta separación es deliberada, ya que PyTorch es una dependencia pesada, pensada para entrenamiento y con un footprint de instalación que no se justifica en un dispositivo que solo hace inferencia, mientras que ONNX Runtime está diseñado específicamente para ese escenario, con un binario liviano y soporte de ejecución en CPU sin necesidad de GPU, coherente con RNF5. Exportar a ONNX tiene además la ventaja de desacoplar el modelo del framework en el que fue entrenado, de modo que si en el futuro conviene migrar la inferencia a otro runtime no hace falta reentrenar ni reexportar CLIP desde cero.

\subsubsection*{Pydantic y SQLite en el contrato de eventos}
El esquema \texttt{MonitoringEvent} que actúa de contrato entre \texttt{vision/}, \texttt{vlm/} y \texttt{baseline/} (Sección~\ref{sec:diagrama_componentes}) está definido con Pydantic, que valida en tiempo de ejecución que cada evento tenga el tipo y la estructura esperada antes de persistirlo, algo que una clase de Python sin validación no garantizaría. La persistencia local usa SQLite a través del módulo estándar \texttt{sqlite3}, sin una capa de ORM intermedia. Se descartó un ORM como SQLAlchemy ya que el volumen y la variedad de consultas sobre \texttt{events/} son bajos (insertar un evento, leerlos todos o filtrados por identificador de video o de cámara), y esas consultas ya están expresadas con claridad en SQL directo, de modo que la capa de abstracción adicional de un ORM no aportaría legibilidad ni velocidad de desarrollo, solo una dependencia más. SQLite en sí, frente a un motor cliente-servidor como PostgreSQL, es la elección natural para este punto del sistema porque el requisito no es concurrencia entre múltiples clientes sino persistencia liviana de un único proceso local (RNF6), sin necesidad de levantar ni administrar un servidor de base de datos en un dispositivo embebido.

\subsubsection*{scikit-learn para el aprendizaje adaptativo}
El módulo \texttt{baseline/}, todavía no iniciado a la fecha de esta entrega (Sección~\ref{sec:diagrama_componentes}), está planificado sobre scikit-learn para implementar el clustering dinámico con ventana móvil descripto en el Marco Teórico (\textcite{Snoun2024}). Una librería de deep learning como PyTorch o TensorFlow resultaría sobredimensionada para esta tarea, dado que el clustering sobre la línea base de comportamiento no requiere entrenar una red neuronal, sino aplicar algoritmos clásicos (por ejemplo DBSCAN o variantes con ventana móvil) sobre un volumen de datos acotado (los eventos diarios de un único residente). Para ese caso de uso, scikit-learn ofrece implementaciones probadas y de bajo costo computacional, coherente con que este módulo también deba poder correr sobre hardware embebido sin GPU.

\subsubsection*{FastAPI, uvicorn y WebSockets en el panel de operación}
El módulo \texttt{dashboard/}, que arranca y detiene \texttt{runner/}, gestiona \texttt{cameras.yaml} y muestra los eventos registrados (RF10), está construido sobre FastAPI con uvicorn como servidor ASGI. Un framework sincrónico como Flask no alcanza acá, porque el panel necesita atender al mismo tiempo peticiones REST convencionales (alta de cámaras, subida de video) y una conexión WebSocket persistente que transmite eventos y logs del runner en tiempo real a la interfaz, y el soporte nativo de FastAPI para operaciones asíncronas es una ventaja directa frente a agregar una extensión externa sobre un framework pensado para trabajar de a una petición por vez. Además, al estar basado en type hints de Python, FastAPI puede reutilizar directamente los modelos Pydantic ya definidos en \texttt{events/}, evitando redefinir la validación de datos en la capa web.

\subsubsection*{Threading y configuración declarativa en la orquestación}
El módulo \texttt{runner/} corre un \textit{thread} por cámara para poder procesar varias fuentes de video en paralelo (RNF9), cada uno con su propia instancia del detector de caídas e inmovilidad, pero compartiendo una única instancia del clasificador CLIP, dado que ONNX Runtime admite inferencia concurrente de forma segura y cargar el modelo una vez por cámara sería redundante. La configuración de cámaras (identificador, fuente, modo de reproducción) se declara en \texttt{cameras.yaml} mediante PyYAML, en lugar de codificarse directamente en el programa, de modo que agregar, quitar o modificar una cámara no requiere cambiar ni redesplegar el código del runner, algo relevante de cara a RF10, que exige poder configurar cámaras desde un panel de administración.

\subsection{Backend en la nube y arquitectura de red}
\label{sec:backend_nube}

Del lado de la nube, el backend que expone \texttt{events/} y \texttt{baseline/} a la aplicación móvil (Sección~\ref{sec:diagrama_arquitectura}) todavía no fue implementado, pero está diseñado sobre FastAPI, el mismo framework que ya corre en \texttt{dashboard/} para el panel de operación local. Reutilizar el mismo framework en ambas puntas tiene una razón concreta, permite compartir el mismo esquema Pydantic de \texttt{MonitoringEvent} entre el backend local y el de la nube sin duplicar la validación, ya que el equipo ya resolvió la curva de aprendizaje de FastAPI a partir de \texttt{dashboard/}, evitando sumar un framework nuevo, por ejemplo Flask o Django, exclusivamente para el backend en la nube sin una razón de peso que lo justifique.

La comunicación entre la Raspberry Pi y el backend en AWS se resolvió sobre REST por HTTPS, y no sobre un protocolo de mensajería como MQTT. La razón es el patrón de tráfico, ya que el sistema no transmite un flujo continuo de datos (el video nunca sale del dispositivo, según RNF1) sino eventos discretos y de baja frecuencia, como una caída, un cambio de actividad o una sincronización periódica de la línea base, un volumen para el que un modelo publish/subscribe agrega complejidad operativa, al requerir un broker adicional para administrar, sin un beneficio claro frente a llamadas HTTP directas. Desde el punto de vista de la arquitectura de red, la Raspberry Pi solo necesita realizar conexiones salientes hacia el endpoint de AWS, sin exponer ningún puerto entrante en la red doméstica, algo especialmente relevante porque el dispositivo corre dentro de la red interna del hogar, típicamente detrás de un router con NAT, sin IP pública ni administración remota de firewall a cargo del usuario. Esa misma conexión saliente es la que se reintenta cuando se restablece la conectividad tras un corte (RNF7), sincronizando los eventos acumulados en SQLite durante ese lapso (RNF6). El tráfico viaja cifrado por TLS, la misma garantía que exige cualquier consumo estándar de una API REST sobre HTTPS, sin necesidad de una capa de cifrado adicional a nivel de aplicación.

Del lado de la persistencia en la nube, el backend va a usar Amazon RDS sobre PostgreSQL en lugar de replicar el SQLite local. La diferencia de requisitos lo justifica, porque en la nube la base de datos tiene que atender accesos concurrentes de la aplicación móvil de distintos hogares al mismo tiempo, algo para lo que SQLite no está pensado, dado que su modelo de concurrencia asume esencialmente un único proceso escritor, mientras que PostgreSQL sí soporta múltiples conexiones concurrentes y permite modelar relaciones entre hogares, usuarios y eventos que en el prototipo local, restringido a un único residente por instalación, todavía no hacen falta.

La autenticación de la aplicación móvil se apoya en Firebase Authentication en lugar de un esquema de tokens propio administrado por el backend, principalmente porque el sistema ya depende de Firebase Cloud Messaging para las notificaciones push (RF8). Sumar también Firebase Authentication evita administrar un segundo proveedor de identidad y aprovecha que el SDK de Firebase en el cliente móvil (React Native) ya integra ambos servicios de forma nativa. El costo es una dependencia mayor de un proveedor externo, pero se trata de un servicio de autenticación estándar, con inicio de sesión por correo y contraseña o proveedores federados, y no de un componente que procese datos sensibles del video ni de la línea base de comportamiento, por lo que no compromete el requisito de privacidad desarrollado en el Marco Teórico (Sección~\ref{sec:edge_marco}).

Por último, la aplicación móvil está desarrollada en React Native. Desarrollar de forma nativa en Swift y en Kotlin por separado no se justifica en este caso, porque la aplicación no requiere acceso a funcionalidades de bajo nivel del dispositivo más allá de recibir notificaciones push y consumir una API REST, y duplicar ese desarrollo tiene poco sentido para un equipo de dos personas con un cronograma acotado. Esta decisión es además independiente de RNF8, que exige que la aplicación sea simple, clara e intuitiva, un objetivo de diseño de interfaz que no depende de si el código corre sobre un framework nativo o multiplataforma.
